\chapter{\DEEN
  {Aufbau wissenschaftlicher Arbeiten}
  {Structure of Academic Papers}}

An der HFH sind verschiedene wissenschaftliche Arbeiten in Form von Haus-,
Projekt- und Abschlussarbeiten zu erbringen. Eine solche wissenschaftliche
Arbeit besteht dabei aus mehreren Teilen. Die verbindliche Reihenfolge der
Teile einer wissenschaftlichen Arbeit an der HFH ist den folgenden
Abschnitten zu entnehmen (Ergänzungen durch Verzeichnisse bzw. Anlagen sind
je nach der von Ihnen behandelten Thematik möglich und zu empfehlen).

\section{\DEEN{Titelblatt}{Title Page}}

Das Titelblatt für themenvereinbarungspflichtige Arbeiten (z.\,B.
Projektarbeit, Abschlussarbeit) erhalten Sie nach Themenanmeldung und
-bestätigung vorgefertigt vom Prüfungsamt. Für andere wissenschaftliche
Arbeiten finden Sie Gestaltungsvorgaben dazu in den separaten Hinweisen zu
den jeweiligen Arbeiten.

Das Titelblatt wird als „Seite 1“ mitgezählt, erhält aber selbst keine
sichtbare Seitenzahl. In dieser LaTeX-Vorlage ersetzt
\texttt{\string\MakeOfficialTitlePage\{datei.pdf\}} das generische Titelblatt
durch die bereitgestellte Seite.

\section{\DEEN{Inhaltsverzeichnis}{Table of Contents}}

Das Inhaltsverzeichnis wird aus den Befehlen
\texttt{\string\chapter}, \texttt{\string\section} und den tieferen
Gliederungsbefehlen automatisch aufgebaut. Nach Änderungen aktualisiert ein
vollständiger \texttt{latexmk}-Lauf die Einträge, Seitenzahlen und Verweise.

Die HFH empfiehlt die Verwendung der Dezimalgliederung -- wie sie auch in
dieser Vorlage genutzt wird. In juristischen Arbeiten wird häufig auch das
alphanumerische System verwendet (siehe dazu Hinweise in den entsprechenden
Studienbriefen zum wissenschaftlichen Arbeiten).

\section{\DEEN
  {Abbildungsverzeichnis und Tabellenverzeichnis}
  {Lists of Figures and Tables}}

Nach dem Inhaltsverzeichnis sind ein Abbildungsverzeichnis und ein
Tabellenverzeichnis eingerichtet. Sie werden automatisch aus
\texttt{\string\caption}-Befehlen erzeugt. Jede Abbildung und Tabelle erhält
zusätzlich ein eindeutiges \texttt{\string\label}; im laufenden Text wird mit
\texttt{\string\ref} darauf verwiesen. Quellen stehen mit
\texttt{\string\source} in einem eigenen Absatz unter dem Objekt und werden
dadurch nicht Bestandteil des Verzeichniseintrags. Nach Änderungen
aktualisiert ein vollständiger \texttt{latexmk}-Lauf beide Verzeichnisse.

\section{\DEEN{Abkürzungsverzeichnis}{List of Abbreviations}}

Eine Erläuterung mit Beispielen finden Sie vorne im
Abkürzungsverzeichnis.

\section{\DEEN
  {Inhaltsteil der wissenschaftlichen Arbeit}
  {Main Body of the Academic Paper}}

Dies ist der wesentliche Teil Ihrer Arbeit, der in der Regel aus einem
einleitenden Teil, einem Hauptteil und einem abschließenden Teil besteht.

\section{\DEEN{Quellenverzeichnis}{References}}

Im Quellenverzeichnis werden alle (und \textit{ausschließlich} die) Quellen
aufgeführt, die Sie in der Arbeit wortwörtlich oder sinngemäß zitiert haben.
Ordnen Sie alle Einträge in das Quellenverzeichnis unabhängig von der Art der
Quelle alphabetisch nach dem Nachnamen des (ersten) Verfassers bzw.
Herausgebers.

Diese Vorlage verwaltet Literatur mit \texttt{biblatex} und Biber. Erfassen
Sie jede Quelle einmal strukturiert in \texttt{references.bib}. Geeignete
Eintragstypen sind etwa Buch (\texttt{@book}), Sammelbandbeitrag
(\texttt{@incollection}), Zeitschriftenaufsatz (\texttt{@article}) und
Onlinequelle (\texttt{@online}). Im Text erzeugt
\texttt{\string\parencite\{schlüssel\}} einen Klammerbeleg;
\texttt{\string\textcite\{schlüssel\}} bindet den Namen in den Satz ein. Ein
vollständiger \texttt{latexmk}-Lauf beziehungsweise Overleaf führt Biber
automatisch aus.

Der Befehl \texttt{\string\printbibliography} erzeugt das
Quellenverzeichnis automatisch an der vorgesehenen Stelle. Es berücksichtigt
standardmäßig nur zitierte Datensätze und sortiert sie nach Name, Jahr und
Titel. Formatieren oder sortieren Sie die Einträge deshalb nicht von Hand.

Der Sammelbeleg
\parencite{booth2016,musterbeitrag2024,musteraufsatz2023,musteronline2025}
demonstriert mehrere Quellen in einem Klammerbeleg. Eine einzelne Quelle
zitieren Sie entweder in Klammern mit \parencite{booth2016} oder im Satz mit
\textcite{booth2016}. Ersetzen Sie die Beispiele in \texttt{references.bib}
durch die tatsächlich verwendeten Quellen. Die Änderungen erscheinen
anschließend automatisch im Quellenverzeichnis.

Die folgenden Muster zeigen die benötigten Angaben für typische Quellentypen:

\begin{itemize}
  \item Name, Vorname (Erscheinungsjahr): Titel der Quelle, ggf. Untertitel.
    Auflage. Erscheinungsort: Verlag.
  \item Name, Vorname (Erscheinungsjahr): Titel des Beitrags/Aufsatzes, ggf.
    Untertitel. In: Name des Herausgebers (Hrsg.): Titel des Sammelwerkes.
    Auflage, Erscheinungsort: Verlag, erste -- letzte Seite des
    Beitrags/Aufsatzes im Sammelwerk.
  \item Name, Vorname (Erscheinungsjahr): Titel des Aufsatzes, ggf. Untertitel.
    In: Titel von Zeitschrift/Journal/Jahrbuch bzw. deren übliche Abkürzung.
    Jahrgang als „Jg.“, ggf. Heftnummer als „H.“, erste -- letzte Seite bzw.
    Spalte des Beitrags.
  \item Name, Vorname (Jahr der Veröffentlichung/Aktualisierung): Titel der
    Internetseite/des Internetdokumentes, ggf. Untertitel. URL: vollständige
    Adresse der konkreten Seite [Stand: Tag Ihres letzten Zugriffs].
\end{itemize}

\section{\DEEN
  {Rechtsprechungsverzeichnis (bei Bedarf)}
  {Table of Cases (if Required)}}

Es enthält eventuell verwendete gerichtliche Entscheidungen (Urteile,
Beschlüsse etc.) mit Fundstelle.

\section{\DEEN
  {Verzeichnis der Verwaltungsanweisungen, Parlamentaria und Ähnliches
   (bei Bedarf)}
  {Administrative and Parliamentary Sources (if Required)}}

Es enthält eventuell verwendete Verwaltungsvorschriften von Landes- oder
Bundesbehörden und eventuell zitierte Parlamentaria (wie Bundestags- oder
Bundesratsdrucksachen).

\section{\DEEN
  {Verzeichnis der eingesetzten KI-Werkzeuge (bei Bedarf)}
  {Register of AI Tools Used (if Required)}}

Wenn Sie KI-Tools als Hilfsmittel bei der Erstellung Ihrer wissenschaftlichen
Arbeit verwenden, sollten Sie deren Einsatz sorgfältig dokumentieren.
Detaillierte Hinweise dazu finden Sie in den „Informationen zum Umgang mit
Künstlicher Intelligenz bei der Erstellung wissenschaftlicher Arbeiten an
der HFH“ -- verfügbar im WebCampus unter den Studienbriefen zum Modul WSA.

\section{\DEEN
  {Anlagenverzeichnis (bei Bedarf)}
  {List of Appendices (if Required)}}

Werden der Arbeit mehrere Anlagen hinzugefügt, sollte ein Anlagenverzeichnis
erstellt werden. Dieses gibt einen Überblick über die der Arbeit beigefügten
Anlagen mit Nummer, Bezeichnung der Anlage und Seitenangabe in systematisch
geordneter Reihenfolge.

\section{\DEEN{Anlagen}{Appendices}}

In den Anlagenteil (Anhang) werden all jene Dokumente aufgenommen, die dabei
helfen, das Verständnis der Arbeit zu steigern, die aber bei einem Einfügen
in den laufenden Text den Gedankenfluss stören würden. Überfrachten Sie den
Anhang allerdings nicht mit Elementen, die anderweitig nicht unterzubringen
sind. Das heißt, alle Anlagen sollten wirklich notwendig sein.

Der Anhang ist nicht dahingehend zu „missbrauchen“, die Seitenbeschränkung
durch Verlagerung von Textpassagen oder Abbildungen in den Anhang zu
unterlaufen.

\section{\DEEN
  {Eigenständigkeitserklärung}
  {Declaration of Authorship}}

Mit der Anfertigung Ihrer wissenschaftlichen Arbeit übernehmen Sie die
Verpflichtung, ehrlich wissenschaftlich zu arbeiten. Plagiieren, d.\,h.
geistiges Eigentum stehlen, gilt als Täuschung und wird gemäß der
Rahmenprüfungsordnung der HFH als nicht bestanden bewertet.

Aus diesem Grund sind Sie verpflichtet, am Ende Ihrer Arbeit eine
Eigenständigkeitserklärung abzugeben.

% ---------------------------------------------------------------------------
% Verzeichnisse nach dem Haupttext
% ---------------------------------------------------------------------------
